\documentclass[11pt]{article}


\usepackage{amsmath,amssymb,amsthm,mathtools}
\usepackage{booktabs,longtable,hyperref,geometry,xcolor,enumitem}
\geometry{margin=1in}
\hypersetup{colorlinks=true,linkcolor=blue,citecolor=blue,urlcolor=blue}
\newtheorem{theorem}{Theorem}[section]
\newtheorem{lemma}[theorem]{Lemma}
\newtheorem{corollary}[theorem]{Corollary}
\newtheorem{proposition}[theorem]{Proposition}
\newtheorem{definition}[theorem]{Definition}
\newtheorem{conjecture}[theorem]{Conjecture}
\theoremstyle{remark}
\newtheorem{remark}[theorem]{Remark}
\newtheorem{observation}[theorem]{Observation}

% Status macros
\newcommand{\LEANVERIFIED}{\textbf{\color{green!60!black}[LEAN-VERIFIED]}}
\newcommand{\PROVED}{\textbf{\color{blue}[PROVED]}}
\newcommand{\NUMERICAL}{\textbf{\color{orange}[NUMERICAL]}}
\newcommand{\HEURISTIC}{\textbf{\color{orange!80!black}[HEURISTIC]}}
\newcommand{\KILLED}{\textbf{\color{red}[KILLED]}}
\newcommand{\OPEN}{\textbf{\color{purple}[OPEN]}}
\newcommand{\CONJECTURE}{\textbf{\color{violet}[CONJECTURE]}}

\title{\textbf{Riemann Hypothesis Research Program:\\
Complete Progress Report and Formal Ledger}}
\author{Logan Eberwein}
\date{September 2026 \\ \small Lean~4 proofs: branch \texttt{stage31-two-channel}, commit \texttt{c579544}}

\begin{document}
\maketitle

\begin{abstract}
This document is a complete research ledger for a multi-year Riemann Hypothesis investigation
combining Lean~4 formal verification, computational search, and analytic mathematics.
All 7 Lean source files (verified) are surveyed; all 190 research markdown documents are synthesized.
Every result is labeled with its verification status:
\LEANVERIFIED{} (machine-checked, no \texttt{sorry}),
\PROVED{} (mathematically proved, not yet formalized),
\NUMERICAL{} (verified by computation),
\HEURISTIC{} (plausible derivation, not rigorous),
\KILLED{} (explicitly disproved with counterexample),
\CONJECTURE{} (stated precisely, status open),
or \OPEN{} (load-bearing open problem).
The Riemann Hypothesis is \OPEN{}.
\end{abstract}

\tableofcontents
\newpage

%%%%%%%%%%%%%%%%%%%%%%%%%%%%%%%%%%%%%%%%%%%%%%%%%%%%%%%%%%
\section{Scope and disclaimer}
%%%%%%%%%%%%%%%%%%%%%%%%%%%%%%%%%%%%%%%%%%%%%%%%%%%%%%%%%%

This is a research report, not a proof of RH.
Numerical positivity at finitely many cutoffs does not imply positivity of the limiting Weil form.
Heuristic asymptotics do not establish analytic theorems.
Every claim is labeled by its exact verification status.
The program has identified the precise mathematical frontier; crossing it requires genuinely new analytic content.

%%%%%%%%%%%%%%%%%%%%%%%%%%%%%%%%%%%%%%%%%%%%%%%%%%%%%%%%%%
\section{Lean-verified theorems: SR/Rees matrix program (Stages 0--40)}
%%%%%%%%%%%%%%%%%%%%%%%%%%%%%%%%%%%%%%%%%%%%%%%%%%%%%%%%%%

The Lean repository (\texttt{D:\textbackslash{}CodexRH}, branch \texttt{stage31-two-channel})
contains 7 Lean source files (verified). The following theorems are machine-verified with no
\texttt{sorry}, \texttt{admit}, or unsafe axiom.

\subsection{Fundamental definitions}

\begin{definition}[Rees comparison matrix] \LEANVERIFIED{}
For $X \geq 6$, define $M_{\mathrm{Rees},X} : \mathrm{Fin}(X-2) \times \mathrm{Fin}(X-2) \to \mathbb{Q}$ by
\[
M_{\mathrm{Rees},X}(i,j) = \begin{cases} -1 & (i+2)(j+2) < X \\ +1 & (i+2)(j+2) \geq X. \end{cases}
\]
\end{definition}

\begin{definition}[Quotient count and threshold set] \LEANVERIFIED{}
\[
\mathrm{SR\_quotient\_count}(X) = \#\bigl\{\lfloor(X-1)/m\rfloor : m \in \{2,\ldots,X-1\}\bigr\}.
\]
The threshold set $\mathrm{SR\_threshold\_set}(X)$ is the image of $m \mapsto \lfloor(X-1)/m\rfloor$
over $\{2,\ldots,X-1\}$.
\end{definition}

\begin{definition}[Threshold submatrix] \LEANVERIFIED{}
For strictly monotone $t_0 < t_1 < \cdots < t_n$:
\[
(S_t)_{ij} = \begin{cases} -1 & t_j \leq t_i \\ +1 & t_j > t_i. \end{cases}
\]
\end{definition}

\begin{definition}[von Mangoldt function] \LEANVERIFIED{}
$\Lambda(n) = \log p$ if $n = p^k$ for some prime $p$ and $k \geq 1$; else $\Lambda(n) = 0$.
\end{definition}

\subsection{Rank theorem}

\begin{theorem}[Rees rank formula] \LEANVERIFIED{} \textnormal{(SR\_rees\_rank\_eq\_quotient\_count)}
For all $X \geq 6$:
\[
\operatorname{rank}(M_{\mathrm{Rees},X}) = \mathrm{SR\_quotient\_count}(X) \approx 2\sqrt{X}.
\]
\end{theorem}

\begin{lemma}[Row equivalence] \LEANVERIFIED{} \textnormal{(SR\_rees\_matrix\_row\_eq\_iff\_quotient\_eq)}
Rows $m$ and $m'$ of $M_{\mathrm{Rees},X}$ are equal iff $\lfloor(X-1)/m\rfloor = \lfloor(X-1)/m'\rfloor$.
\end{lemma}

\begin{lemma}[Threshold row difference] \LEANVERIFIED{} \textnormal{(SR\_rees\_threshold\_row\_diff)}
For quotients $q_1 < q_2$ and index $j$:
\[
R_{q_2}(j) - R_{q_1}(j) = \begin{cases} -2 & q_1 < j+2 \leq q_2 \\ 0 & \text{otherwise.} \end{cases}
\]
\end{lemma}

\begin{theorem}[Quotient rows are linearly independent] \LEANVERIFIED{} \textnormal{(SR\_rees\_quotient\_rows\_linearIndependent)}
The threshold representative rows of $M_{\mathrm{Rees},X}$ are linearly independent over $\mathbb{Q}$.
\end{theorem}

\subsection{Threshold matrix: coercivity theorems}

\begin{theorem}[Symmetric part identity] \LEANVERIFIED{} \textnormal{(SR\_threshold\_matrix\_sym\_add\_self)}
For any strictly monotone $t : \mathrm{Fin}(n+1) \to \mathbb{N}$:
\[
S_t + S_t^T = -2I.
\]
\end{theorem}

\begin{theorem}[Exact quadratic form] \LEANVERIFIED{} \textnormal{(SR\_threshold\_matrix\_quadratic\_form)}
\[
x^T S_t x = -\|x\|^2 \quad \text{for all } x \in \mathbb{Q}^{n+1}.
\]
\end{theorem}

\begin{theorem}[Trivial kernel / injectivity] \LEANVERIFIED{} \textnormal{(SR\_threshold\_matrix\_mulVec\_injective)}
$S_t$ is injective: $S_t x = 0 \Rightarrow x = 0$.
\end{theorem}

\begin{theorem}[Skew decomposition] \LEANVERIFIED{} \textnormal{(SR\_threshold\_matrix\_eq\_neg\_one\_add\_skew)}
$S_t = -I + K$ where $K = (S_t - S_t^T)/2$ is skew-symmetric ($K^T = -K$).
\end{theorem}

\begin{corollary} \LEANVERIFIED{}
The skew part satisfies $K^T = -K$ (SR\_threshold\_matrix\_skew\_transpose).
\end{corollary}

\begin{theorem}[Signed threshold determinant] \LEANVERIFIED{} \textnormal{(SR\_signed\_threshold\_matrix\_det)}
\[
\det(S_t) = (-1) \cdot (-2)^n \neq 0.
\]
\end{theorem}

\begin{theorem}[Linear independence from determinant] \LEANVERIFIED{} \textnormal{(SR\_signed\_threshold\_matrix\_linearIndependent)}
The rows of $S_t$ are linearly independent over $\mathbb{Q}$.
\end{theorem}

\subsection{Global positivity}

\begin{theorem}[All-ones positivity] \LEANVERIFIED{} \textnormal{(allOnesEntrySum\_positive\_general)}
For all $X \geq 6$:
\[
A_X := \sum_{m,n \in \{2,\ldots,X-1\}} M_{\mathrm{Rees},X}(m,n) > 0.
\]
\end{theorem}

\begin{observation} \NUMERICAL{}
Machine-checked: $A_X > 0$ for all $6 \leq X \leq 1000$.
Benchmark values: $A_6 = 5$, $A_{10} = 35$, $A_{30} = 635$, $A_{100} = 8855$,
$A_{210} = 41381$, $A_{2310} = 5295017$.
\end{observation}

\begin{theorem}[Hyperbola reduction] \LEANVERIFIED{}
$A_X = (X-2)^2 - 2N_{\mathrm{int}}(X)$ where $N_{\mathrm{int}}(X) = \#\{(m,n): mn < X, m,n \in \mathrm{support}\}$.
The positivity follows from $N_{\mathrm{int}}(X) < (X-2)^2/2$ via the Dirichlet hyperbola method.
\end{theorem}

\subsection{Indefiniteness: impossibility theorems}

These are among the most important results: they rule out entire classes of proof strategies.

\begin{theorem}[Negative direction exists] \LEANVERIFIED{} \textnormal{(SR\_form\_has\_negative\_direction)}
For all $X \geq 6$, there exists $v : \mathrm{Fin}(X-2) \to \mathbb{R}$ such that
\[
\operatorname{SR\_quadratic\_form}(X, v) < 0.
\]
The basis vector $e_0$ (supported at index $0$, which corresponds to $m=2$) witnesses this.
\end{theorem}

\begin{theorem}[Raw SR form is not PSD] \LEANVERIFIED{} \textnormal{(SR\_form\_not\_positive\_semidefinite)}
For all $X \geq 6$:
\[
\neg\,\bigl(\forall v,\; 0 \leq \operatorname{SR\_quadratic\_form}(X, v)\bigr).
\]
\end{theorem}

\begin{theorem}[No PSD realization preserving raw form] \LEANVERIFIED{} \textnormal{(SR\_no\_psd\_realization\_preserving\_raw\_form)}
Let $X \geq 6$. There is no positive semidefinite form $P$ satisfying
$P(v) = \operatorname{SR\_quadratic\_form}(X, v)$ for all $v$.
Formally: for any $P : (\mathrm{Fin}(X-2) \to \mathbb{R}) \to \mathbb{R}$ with $P(v) \geq 0$ everywhere,
$P \neq \operatorname{SR\_quadratic\_form}(X, \cdot)$.
\end{theorem}

\begin{theorem}[Global coercivity is false] \LEANVERIFIED{} \textnormal{(SR\_global\_coercivity\_false)}
\[
\neg\,\bigl(\forall X \geq 6,\; \forall v,\; 0 \leq \operatorname{SR\_quadratic\_form}(X, v)\bigr).
\]
\end{theorem}

\begin{theorem}[SR form is indefinite] \LEANVERIFIED{} \textnormal{(SR\_form\_is\_indefinite\_of\_positive)}
For all $X \geq 6$: whenever the SR quadratic form takes a positive value,
it also takes a negative value. The form is indefinite.
\end{theorem}

\begin{theorem}[Explicit $X=6$ witness] \LEANVERIFIED{} \textnormal{(SR\_X6\_indefinite)}
At $X = 6$: $\operatorname{SR\_quadratic\_form}(6, e_0) < 0$ and $\operatorname{SR\_quadratic\_form}(6, \mathbf{1}) = 14 > 0$.
\end{theorem}

\begin{theorem}[Uniform subspace is positive] \LEANVERIFIED{} \textnormal{(SR\_uniform\_X6\_positive, SR\_uniform\_X6\_strict)}
For $X=6$, on the 1-dimensional subspace of constant vectors:
$\operatorname{SR\_quadratic\_form}(6, c\cdot\mathbf{1}) = 14c^2 \geq 0$ with equality iff $c=0$.
\end{theorem}

\subsection{Dagger involution (bidegree analysis)}

\begin{theorem}[Dagger is an involution] \LEANVERIFIED{} \textnormal{(dagger\_involution)}
On the $(1,1)$ carrier $E^{11}$ at $X=6$: $\dagger(\dagger(x)) = x$ for all $x$.
\end{theorem}

\begin{theorem}[Dagger swaps labels] \LEANVERIFIED{} \textnormal{(dagger\_h\_swap)}
For all $m, n$: $\dagger(h_{m\bar{n}}) = h_{n\bar{m}}$.
\end{theorem}

\begin{theorem}[Diagonal elements are fixed] \LEANVERIFIED{} \textnormal{(dagger\_diagonal\_fixed)}
For all $m$: $\dagger(h_{m\bar{m}}) = h_{m\bar{m}}$.
\end{theorem}

\begin{theorem}[Dagger is $\mathbb{C}$-antilinear] \LEANVERIFIED{}
$\dagger(x + y) = \dagger(x) + \dagger(y)$ and $\dagger(a \cdot x) = \bar{a} \cdot \dagger(x)$.
\end{theorem}

\subsection{GNS-type positivity}

\begin{theorem}[GNS positive mass] \LEANVERIFIED{} \textnormal{(GNS\_positive)}
For all $x \in E^{11}$: $0 \leq |\Gamma_{\mathrm{obs}}(x)|^2$.
\end{theorem}

\begin{theorem}[GNS dagger-positivity] \LEANVERIFIED{} \textnormal{(GNS\_positive\_correct)}
For all $x \in E^{11}$: $0 \leq \operatorname{Re}(\Gamma_{\mathrm{obs}}(\dagger(x) \cdot x))$.
\end{theorem}

\begin{theorem}[Type-level firewall] \LEANVERIFIED{} \textnormal{(Gamma\_obs\_domain\_E11)}
The GNS functional $\Gamma_{\mathrm{obs}}$ is defined on $E^{11}$, not on $E^{20}_6$.
This is a type-level obstruction: the $(1,1)$-sector extension principle is not automatic.
\end{theorem}

\subsection{Zero-transfer (margin argument)}

\begin{theorem}[Nonzero transfer] \LEANVERIFIED{} \textnormal{(nonzero\_transfer\_of\_margin)}
In any normed ring: if $\delta > 0$, $\delta \leq \|F\|$, and $\|G - F\| < \delta$,
then $G \neq 0$.
\end{theorem}

\begin{theorem}[Uniform zero-free transfer] \LEANVERIFIED{} \textnormal{(uniform\_zero\_free\_transfer)}
Pointwise: if $\delta \leq \|F(x)\|$ and $\|G(x) - F(x)\| < \delta$ for all $x$,
then $G(x) \neq 0$ for all $x$.
\end{theorem}

\subsection{Arithmetic bridge and channel decompositions}

\begin{definition}[SR Sint and Sext] \LEANVERIFIED{}
\begin{align*}
S_{\mathrm{int},\log}(X) &= \sum_{\substack{m,n \in \mathrm{supp}(X) \\ mn < X}} \frac{\log m \cdot \log n}{\sqrt{mn}}, \\
S_{\mathrm{ext},\log}(X) &= \sum_{\substack{m,n \in \mathrm{supp}(X) \\ mn \geq X}} \frac{\log m \cdot \log n}{\sqrt{mn}}.
\end{align*}
\end{definition}

\begin{theorem}[B3 bound] \LEANVERIFIED{} \textnormal{(B3\_final, Stage 27)}
\[
\frac{S_{\mathrm{int},\log}(X)}{T_{\log}(X)} \leq \frac{1 - \sqrt{\kappa_{\log}(X)}}{2}.
\]
\end{theorem}

\begin{theorem}[Kappa bounds] \LEANVERIFIED{}
$\kappa_{\log}(X) \in [0, 1]$ for all $X$.
\end{theorem}

\begin{theorem}[Two-channel decomposition] \LEANVERIFIED{} \textnormal{(SR\_two\_channel\_decomposition, Stage 28)}
\[
Q = \tfrac{1}{2}Q_{\mathrm{ratio}} + \tfrac{1}{2}Q_{\mathrm{product}}
\]
via the identity $\cos A \cos B = \frac{1}{2}\cos(A-B) + \frac{1}{2}\cos(A+B)$.
\end{theorem}

\begin{theorem}[Ratio channel diagonal bound] \LEANVERIFIED{} \textnormal{(SR\_ratio\_diagonal\_bound, Stage 33)}
The diagonal of $Q_{\mathrm{ratio}}$ satisfies:
\[
\mathrm{diag}(Q_{\mathrm{ratio}})/X^{2\beta} \to \frac{1}{2\beta} \quad \text{as } X \to \infty.
\]
\end{theorem}

\subsection{Dirichlet series identities}

\begin{theorem}[von Mangoldt L-series] \LEANVERIFIED{} \textnormal{(sr37\_clean\_vonMangoldt\_lseries\_eq\_negative\_zeta\_log\_deriv)}
In the half-plane $\operatorname{Re}(s) > 1$:
\[
\sum_{n=1}^\infty \frac{\Lambda(n)}{n^s} = -\frac{\zeta'(s)}{\zeta(s)}.
\]
\end{theorem}

\begin{theorem}[Convolution L-series] \LEANVERIFIED{} \textnormal{(sr37\_clean\_unrestricted\_pair\_convolution\_lseries\_eq\_zeta\_sq)}
\[
L(\mathbf{1} * \mathbf{1}, s) = \zeta(s)^2.
\]
\end{theorem}

\begin{theorem}[Finite Mellin kernels] \LEANVERIFIED{} \textnormal{(Stage 36)}
Sharp-cutoff finite Mellin transform $\to 1/s$.
First-difference finite Mellin transform $\to 1/(s(s+1))$.
Second-difference finite Mellin transform $\to 2/(s(s+1)(s+2))$.
\end{theorem}

\begin{theorem}[SR Dirichlet series limit] \PROVED{}
As $X \to \infty$, the SR Dirichlet series converges to:
\[
D_{\mathrm{SR}}(s) \to -(\zeta(s) - 1)^2 \quad \text{(support starts at } m=2\text{)}.
\]
With full support $m \geq 1$: $D_{\mathrm{full}}(s) \to -\zeta(s)^2$.
\end{theorem}

\begin{theorem}[Signature data at primorials] \NUMERICAL{}
The Rees matrix signature $(p, q, r)$ is balanced ($p = q$) at primorial values
$X = 6, 30, 210, 2310$ under complete integer support.
Verified: $X = 6, 30, 210, 2310$ all give $p = q$.
Non-primorials $X = 7, 35$ give $p \neq q$.
\end{theorem}

%%%%%%%%%%%%%%%%%%%%%%%%%%%%%%%%%%%%%%%%%%%%%%%%%%%%%%%%%%
\section{Proved results: analytic program}
%%%%%%%%%%%%%%%%%%%%%%%%%%%%%%%%%%%%%%%%%%%%%%%%%%%%%%%%%%

\subsection{Weil form contraction search (25+ candidates)}

\begin{theorem}[Hardy-Gram exact multiplicativity] \PROVED{}
For all positive integers $c, m, n$:
\[
G_H(cm, cn) = c \cdot G_H(m, n),
\]
where
\[
G_H(m,n) = \frac{\sqrt{mn}}{2\pi} \int_{-\infty}^\infty \left(\frac{m}{n}\right)^{it}
\frac{|\zeta(\tfrac{1}{2}+it)|^2}{\tfrac{1}{4}+t^2}\,dt.
\]
Proof: $\sqrt{cm \cdot cn} = c\sqrt{mn}$ and $(cm/cn)^{it} = (m/n)^{it}$. $\square$
\end{theorem}

\begin{corollary} \PROVED{}
$G_H(m,n) = \gcd(m,n) \cdot G_H(m/\gcd, n/\gcd)$.
\end{corollary}

\begin{theorem}[Hardy-Gram structural form] \PROVED{}
\[
G_H(m,n) = \sqrt{mn} \cdot H\!\left(\log\frac{m}{n}\right), \quad
H(\tau) = \frac{1}{2\pi}\int_{-\infty}^\infty \frac{|\zeta(\tfrac{1}{2}+it)|^2}{\tfrac{1}{4}+t^2} e^{i\tau t}\,dt.
\]
$H$ is the Fourier transform of $|\zeta(1/2+it)|^2/(1/4+t^2)$.
\end{theorem}

\begin{theorem}[Hardy-Gram asymptotics] \HEURISTIC{}
From the approximate functional equation: as the integration cutoff $T \to \infty$,
\[
\frac{G_H(m,n)}{G_H(1,1)} \to \gcd(m,n).
\]
The normalized Hardy Gram satisfies $G_H(m,n)/\sqrt{G_H(m,m) G_H(n,n)} \to \gcd(m,n)/\sqrt{mn}$.
\end{theorem}

\begin{observation} \NUMERICAL{}
Diagonal: $G_H(n,n)/G_H(1,1) = n$ exactly (by multiplicativity).
Off-diagonal: $G_H(1,2)/G_H(1,1)$ decreases from $1.230$ at $T=100$ to $1.204$ at $T=1000$,
converging toward $\gcd(1,2) = 1$.
\end{observation}

\begin{theorem}[GCD kernel failure margin] \NUMERICAL{}
The Weil form contraction norm with the GCD kernel satisfies:
\[
\|C_L\| - 1 \sim L^{7.30} \quad (L \in [6.5, 8.0], \text{ 4-point power-law fit}).
\]
\end{theorem}

\begin{theorem}[Dominant failure modes at small primes] \NUMERICAL{}
The dominant negative eigenvector of $G_L^* G_L - F_L^* F_L$ is concentrated on
prime-power channels with $n = 2, 4, 7, 9, 13, 16, 17, 19$
consistently across $L = 6, 7, 8$. These are NOT boundary primes near $\exp(2L)$.
\end{theorem}

\subsection{Killed approaches: exact failure theorems}

\begin{theorem}[SR coefficients $\neq \Lambda$] \PROVED{} (counterexample)
$\operatorname{SR\_coeff}(6, 4) = -1 \neq \Lambda(4) = \log 2$.
SR produces divisor-count coefficients, not von Mangoldt weights.
\end{theorem}

\begin{theorem}[Abel bridge wrong scale] \PROVED{} (Stage 39 disproof)
$S_{\mathrm{int},\log}(X) \sim 2\sqrt{X}\log X$, while $2\psi_2(X)/\sqrt{X} \sim X^{3/2}$.
These differ by a factor of $X$. The Abel bridge approximation is false.
\end{theorem}

\begin{theorem}[$\zeta^2$ gives zeros not poles at $\zeta$ zeros] \PROVED{}
$D_{\mathrm{full}}(s) = -\zeta(s)^2$ has \emph{zeros} of order 2 at $\zeta$ zeros,
not poles. Perron/explicit formula gives no contribution from $\zeta$ zeros.
\end{theorem}

\begin{theorem}[Weil/PW clamped inequality is false for small $L$] \PROVED{}
For $h(y) = y(1-y^2)^2$: as $L \to 0$, $R_L \to 337/105 < \infty$
while $c_L \sim \pi^2/(2L) \to \infty$.
For all $0 < L \leq 1/10$: $R_L < 4$ while $c_L > 20$.
The inequality $E_L[HF] \geq c_L\|HF\|^2$ is false for all sufficiently small $L$.
\end{theorem}

\begin{theorem}[Relaxed-source scaling obstruction] \PROVED{} (Research Ledger)
For $D = ch$ and any positive scaling $\lambda > 0$: the LP inequality
$\lambda^2(e - \ell^2) + 8\lambda\ell - 16 \leq 0$ fails for large positive $\lambda$ whenever $e > \ell^2$.
Any theorem extending contractivity to arbitrary scaled sources is false.
\end{theorem}

\begin{theorem}[Volterra endpoint mismatch] \PROVED{} (Research Ledger)
The causal convolution representation $V_\alpha$ satisfies $V_{\delta_T} = 0$ as an operator on $L^2(0,T)$.
Hence $V_{\nu_M} = V_{\mu_M}$ despite $\cos$-channel difference $V_M(0) = C_M(0) - m_M \cos(0) \neq 0$.
The Volterra representation is non-faithful on terminal boundary mass.
Smallest counterexample: $\alpha = \delta_T$ with $0 < M \leq (\log 2)/2$.
\end{theorem}

\begin{theorem}[Finite moment closure killed] \PROVED{} (Research Ledger)
A polynomial identity in the endpoint-centered moments $b_k = \int (T-s)^k d\mu(s)$
that closes the terminal ladder at any finite order would require a polynomial relation
vanishing on the continuum support interval --- hence the zero polynomial.
Therefore the exact linear scale/endpoint state cannot be finite-dimensional.
\end{theorem}

\begin{theorem}[Ratio-channel GUE detection is $\gamma$-independent] \PROVED{}
The matrix $M_{\mathrm{Rees}}(m,n) \cdot e^{i\gamma(\log m - \log n)}$ is unitarily equivalent
(via diagonal conjugation) to $M_{\mathrm{Rees}}$ for every $\gamma$.
Its spectrum, and hence its $r$-statistic, is $\gamma$-independent.
\end{theorem}

\begin{theorem}[Product channel non-specificity] \NUMERICAL{}
The product-channel kernel $\mathrm{Rees}(m,n) \cdot \cos(\gamma \log(mn)/2)$:
at $X=200$, $\gamma = 35.0$ (not a $\zeta$ zero) gives $r = 0.6054$,
closer to GUE ($= 0.5996$) than any tested $\zeta$ zero frequency.
\end{theorem}

\begin{theorem}[M\"obius-log operator is indefinite] \PROVED{}
Any operator with trace $= \sum_k \Lambda(k)/k^s$ must be indefinite.
Diagonal entries involve $\mu(m) \log m$ which takes both signs:
$\mu(2)\log 2 = -\log 2 < 0$ but $\mu(30)\log 30 = -\log 30 < 0$.
\end{theorem}

\begin{theorem}[BN target vector is tautological in gcd model] \PROVED{}
In the gcd kernel model with $b_k = 1/\sqrt{k}$:
$b = $ first column of $G$, so $G^{-1}b = e_1$ and $b^T G^{-1} b = G_{11} = 1$.
Hence $d_N = 0$ for all $N$ by linear algebra alone.
This is not evidence of BN convergence.
\end{theorem}

\begin{theorem}[Arakelov matrix is indefinite] \NUMERICAL{}
The matrix $\mathrm{Ar}(m,n) = \frac{1}{2}\log(\mathrm{lcm}(m,n)/\gcd(m,n)^2) \cdot \gcd(m,n)/\sqrt{mn}$
has minimum eigenvalue $\approx -5.49$ at $N=100$ with $99$ negative eigenvalues.
It is strongly indefinite.
\end{theorem}

\subsection{Structural obstruction theorems}

\begin{theorem}[Type mismatch] \PROVED{}
RH is a \emph{support condition}: zeros lie on $\operatorname{Re}(s) = 1/2$.
Matrix positivity is a \emph{sign condition}: eigenvalues $\geq 0$.
These require fundamentally different mathematical machinery:
support conditions need global analytic continuation;
sign conditions are local and finite.
The claim is narrow: this specific SR comparison matrix has no currently proved
mechanism converting its finite sign information into the required analytic support statement.
Note that Li's criterion is a positivity condition equivalent to RH, and the Redheffer matrix
determinant equals the Mertens function---so finite matrices can encode RH-equivalent information.
\end{theorem}

\begin{theorem}[Separability obstruction] \PROVED{}
$\Lambda(mn) \neq \Lambda(m)\Lambda(n)$ for general $m,n$.
No separable kernel $f(m)g(n)$ applied to any matrix with entries depending
separately on $m$ and $n$ can produce $\Lambda$ as a coefficient sequence.
\end{theorem}

\begin{theorem}[Circularity obstruction] \PROVED{}
The Beurling-Nyman target $b_n = \langle 1_{[0,1]}, \rho_n\rangle_{H^2}$ requires
evaluating $\zeta(1/2+it)$ on the critical line.
Computing this without resolving the zero locations is equivalent to knowing RH.
The target cannot be computed from first principles without circular assumption.
\end{theorem}

\begin{theorem}[Arithmetic/spectral duality obstruction] \PROVED{} (Research Ledger)
The Weil explicit formula's arithmetic side (primes, poles, gamma) and spectral side (zeros)
are dual descriptions of the same functional. They are not independent positive-energy terms
that may be added. Any attempt to supplement the arithmetic form with zero-contribution
energy changes the form rather than improving the original one.
\end{theorem}

%%%%%%%%%%%%%%%%%%%%%%%%%%%%%%%%%%%%%%%%%%%%%%%%%%%%%%%%%%
\section{The required operator: formal criteria}
%%%%%%%%%%%%%%%%%%%%%%%%%%%%%%%%%%%%%%%%%%%%%%%%%%%%%%%%%%

\begin{definition}[Six required properties for the correct operator]
An operator $T(s)$ could prove RH if it satisfies all of:

\begin{enumerate}[label=\textbf{P\arabic*}]
\item \textbf{[Poles at zeros]} $\det(I-T(s)) = 0$ at each $\zeta$ zero $\rho$.
\item \textbf{[Prime-power weights]} $\operatorname{Tr}(T^n(s))$ counts prime powers with weight $\Lambda(p^k)$.
\item \textbf{[Positivity]} Some canonical quadratic form $Q[T] \geq 0$.
\item \textbf{[Functional equation]} $T(s)$ and $T(1-s)$ are related; $\operatorname{Re}(s)=1/2$ is the fixed line.
\item \textbf{[Arithmetic origin]} $T$ is constructible from integers; finite truncations $T_X \to T$.
\item \textbf{[Convergence]} $\det(I-T(s))$ is well-defined and analytically continuable through the critical strip.
\end{enumerate}
\end{definition}

\begin{theorem}[Property table] \NUMERICAL{}/\PROVED{}
\begin{center}
\begin{tabular}{lcccccc}
\toprule
Operator & P1 & P2 & P3 & P4 & P5 & P6 \\
\midrule
SR Rees $M_{\mathrm{Rees},X}$ & \texttimes & \texttimes & $\sim$ & \texttimes & \checkmark & \texttimes \\
Full-support SR & \texttimes & \texttimes & $\sim$ & \texttimes & \checkmark & \texttimes \\
Threshold submatrix $S_t$ & \texttimes & \texttimes & \checkmark & \texttimes & \checkmark & \texttimes \\
Weil/PW $H=-D^2+1/4$ & $\sim$ & \texttimes & $\sim$ & \checkmark & \texttimes & \checkmark \\
Mayer transfer $\mathcal{L}_s$ & \checkmark & \checkmark & \texttimes & \checkmark & \texttimes & \checkmark \\
Connes adelic $L^2(A_\mathbb{Q}/\mathbb{Q}^\times)$ & \checkmark & \checkmark & ? & \checkmark & \texttimes & \checkmark \\
GCD matrix $\gcd(m,n)/\sqrt{mn}$ & $\sim$ & \checkmark & \checkmark & \texttimes & \checkmark & \texttimes \\
Hardy-Gram $G_H$ (normalized) & \checkmark & $\sim$ & \checkmark & $\sim$ & $\sim$ & $\sim$ \\
\bottomrule
\end{tabular}
\end{center}
\textbf{No existing operator satisfies all six properties.}
\end{theorem}

%%%%%%%%%%%%%%%%%%%%%%%%%%%%%%%%%%%%%%%%%%%%%%%%%%%%%%%%%%
\section{Conjectures: precisely stated open problems}
%%%%%%%%%%%%%%%%%%%%%%%%%%%%%%%%%%%%%%%%%%%%%%%%%%%%%%%%%%

\begin{conjecture}[Primorial balance] \CONJECTURE{}
The Rees matrix signature satisfies $p = q$ (balanced) at every primorial $X = P_k = 2 \cdot 3 \cdot 5 \cdots p_k$
under complete integer support $\{1, \ldots, X-1\}$.
Verified: $k = 2, 3, 4$ ($X = 6, 30, 210, 2310$). Open for all $k$.
\end{conjecture}

\begin{conjecture}[Universal signature bound] \CONJECTURE{}
For all $X \geq 6$: $|\text{pos}(X) - \text{neg}(X)| \leq 1$,
where pos and neg count positive and negative eigenvalues of $M_{\mathrm{Rees},X}$.
Verified numerically for $X \leq 1000$.
\end{conjecture}

\begin{conjecture}[Dagger-completion defect] \CONJECTURE{}
The fundamental obstruction across all major RH reformulations is the
\emph{dagger-completion defect}: the inability to convert forward multiplicative
structure ($A^2$) into Hermitian/dagger structure ($A^\dagger A$).
This gap, denoted $\mathfrak{D}_\dagger(A)$, is conjectured to be the universal failure invariant.
Status: identified across SR, Connes, Weil programs. Not yet formalized.
\end{conjecture}

\begin{conjecture}[Hardy-Gram GCD asymptotics] \CONJECTURE{}
\[
\frac{G_H(m,n)}{G_H(1,1)} \to \gcd(m,n) \quad \text{as } T \to \infty,
\]
rigorously, from the asymptotic theory of the twisted second moment of $\zeta$.
The heuristic derivation from the approximate functional equation is given in \S3.
A rigorous proof would follow from the Montgomery-Vaughan mean value theorem
for the twisted second moment.
\end{conjecture}

%%%%%%%%%%%%%%%%%%%%%%%%%%%%%%%%%%%%%%%%%%%%%%%%%%%%%%%%%%
\section{The frontier}
%%%%%%%%%%%%%%%%%%%%%%%%%%%%%%%%%%%%%%%%%%%%%%%%%%%%%%%%%%

The following is the load-bearing open equation from the research ledger.

\begin{theorem}[Load-bearing open equation] \OPEN{} (Research Ledger)
Construct from the completed arithmetic source a contractive system
whose first derivative maps a structurally derived input $b$ to
\[
x_{\mathrm{obs}} = \bigl(\sqrt{B}\,D,\; a \mapsto \sqrt{B}\,P(\Theta_a V_M)\bigr)
\]
and whose independently derived input norm is exactly $\|b\| = |4 - L|$.
A generic Schur factorization with these properties is merely equivalent to LP
and is not evidence for it.
\end{theorem}

\textbf{The three-step frontier:}

\begin{enumerate}
\item \textbf{Step 1} (Requires audit first): Certify or kill the Hardy-Gram asymptotics (Conjecture 4.4).
The conjecture $G_H(m,n)/G_H(1,1) \to \gcd(m,n)$ is not yet supported by convergent numerics:
the ratio $G_H(1,2)/G_H(1,1) \approx 1.204$ at $T=1000$ remains far from the conjectured limit of $1$.
Resolve before treating this as a connection between the Weil contraction and BN programs.

\item \textbf{Step 2} (Open): Compute the correct BN target vector $b_n = \langle 1_{[0,1]}, \rho_n\rangle_{H^2}$
analytically, without assuming RH.

\item \textbf{Step 3} (= RH): Prove $d_N \to 0$, i.e., the Beurling-Nyman approximation converges.
This is equivalent to RH by the Báez-Duarte criterion.
The gap between Steps 2 and 3 is the Riemann Hypothesis itself.
\end{enumerate}

%%%%%%%%%%%%%%%%%%%%%%%%%%%%%%%%%%%%%%%%%%%%%%%%%%%%%%%%%%
\section{Repository and reproducibility}
%%%%%%%%%%%%%%%%%%%%%%%%%%%%%%%%%%%%%%%%%%%%%%%%%%%%%%%%%%

\begin{itemize}
\item \textbf{Repository:} \url{https://github.com/loganeberwein6/Logan-Eberwein-RH-Research}
\item \textbf{Branch:} \texttt{stage31-two-channel}, commit \texttt{c579544}
\item \textbf{Lean files:} 7 (verified, no \texttt{sorry})
\item \textbf{Documented approaches:} 190 markdown files in \texttt{approaches/}
\item \textbf{Build:} Set \texttt{ELAN\_HOME} to your elan installation path, then \texttt{lake update} and \texttt{lake build}.
\item \textbf{Zenodo DOI:} \texttt{10.5281/zenodo.22783121}
\end{itemize}

%%%%%%%%%%%%%%%%%%%%%%%%%%%%%%%%%%%%%%%%%%%%%%%%%%%%%%%%%%
\section{Summary table}
%%%%%%%%%%%%%%%%%%%%%%%%%%%%%%%%%%%%%%%%%%%%%%%%%%%%%%%%%%

\begin{longtable}{p{8cm}ll}
\toprule
Result & Status & Stage \\
\midrule
Rees matrix definition and basic properties & \LEANVERIFIED{} & 0--10 \\
$\operatorname{rank}(M_{\mathrm{Rees},X}) = \mathrm{SR\_quotient\_count}(X)$ & \LEANVERIFIED{} & 38 \\
$S_t + S_t^T = -2I$ & \LEANVERIFIED{} & 38 \\
$x^T S_t x = -\|x\|^2$ & \LEANVERIFIED{} & 38 \\
$S_t$ injective & \LEANVERIFIED{} & 38 \\
$\det(S_t) = (-1)(-2)^n$ & \LEANVERIFIED{} & 38 \\
Threshold rows linearly independent & \LEANVERIFIED{} & 38 \\
$A_X > 0$ for all $X \geq 6$ & \LEANVERIFIED{} & 21--23 \\
$A_X > 0$ machine-checked to $X=1000$ & \NUMERICAL{} & 23 \\
Raw SR form is not PSD & \LEANVERIFIED{} & 32 \\
No PSD realization preserves raw SR form & \LEANVERIFIED{} & 32 \\
Global coercivity is false & \LEANVERIFIED{} & 32 \\
SR form is indefinite & \LEANVERIFIED{} & 32 \\
$X=6$ explicit indefiniteness witness & \LEANVERIFIED{} & 32 \\
Dagger involution on $E^{11}$ & \LEANVERIFIED{} & 39--40 \\
GNS dagger-positivity & \LEANVERIFIED{} & 39--40 \\
Type-level domain firewall & \LEANVERIFIED{} & 40 \\
Zero-transfer margin theorem & \LEANVERIFIED{} & 40 \\
$L(\Lambda,s) = -\zeta'/\zeta$ & \LEANVERIFIED{} & 37 \\
$L(\mathbf{1}*\mathbf{1},s) = \zeta(s)^2$ & \LEANVERIFIED{} & 37 \\
B3 bound: $S_{\mathrm{int}}/T \leq (1-\sqrt{\kappa})/2$ & \LEANVERIFIED{} & 27 \\
$\kappa \in [0,1]$ & \LEANVERIFIED{} & 27 \\
Two-channel decomposition & \LEANVERIFIED{} & 28 \\
Finite Mellin kernels & \LEANVERIFIED{} & 36 \\
Primorial balance at $X=6,30,210,2310$ & \NUMERICAL{} & 19--23 \\
Signature bound $|\text{pos}-\text{neg}| \leq 1$ ($X \leq 1000$) & \NUMERICAL{} & 17--23 \\
\midrule
$G_H(cm,cn) = c \cdot G_H(m,n)$ & \PROVED{} & This work \\
$G_H(m,n) = \sqrt{mn}\,H(\log(m/n))$ & \PROVED{} & This work \\
$G_H(m,n)/G_H(1,1) \to \gcd(m,n)$ & \HEURISTIC{} & This work \\
Weil failure margin $\sim L^{7.3}$ & \NUMERICAL{} & This work \\
Small-prime failure mode identification & \NUMERICAL{} & This work \\
Weil/PW clamped inequality false for small $L$ & \PROVED{} & This work \\
Abel bridge wrong scale & \PROVED{} & Stage 39 \\
$\zeta^2$ gives zeros not poles & \PROVED{} & This work \\
Ratio-channel GUE $\gamma$-independence & \PROVED{} & This work \\
M\"obius-log operator indefinite & \PROVED{} & This work \\
BN gcd model is tautological & \PROVED{} & This work \\
Type mismatch obstruction & \PROVED{} & This work \\
Separability obstruction & \PROVED{} & This work \\
Circularity obstruction & \PROVED{} & This work \\
Volterra endpoint mismatch & \PROVED{} & Research Ledger \\
Finite moment closure impossible & \PROVED{} & Research Ledger \\
Relaxed-source scaling failure & \PROVED{} & Research Ledger \\
Arakelov matrix indefinite & \NUMERICAL{} & This work \\
\midrule
Dagger-completion defect conjecture & \CONJECTURE{} & Ledger \\
Primorial balance conjecture & \CONJECTURE{} & Stage 19 \\
Universal signature bound conjecture & \CONJECTURE{} & Stage 17 \\
Hardy-Gram GCD asymptotics & \CONJECTURE{} & This work \\
\midrule
RH & \OPEN{} & --- \\
\bottomrule
\end{longtable}

%%%%%%%%%%%%%%%%%%%%%%%%%%%%%%%%%%%%%%%%%%%%%%%%%%%%%%%%%%
\section*{Acknowledgments}
%%%%%%%%%%%%%%%%%%%%%%%%%%%%%%%%%%%%%%%%%%%%%%%%%%%%%%%%%%

Mathematical development: Logan Eberwein.
Lean formalization: automated via GPT-4o and Claude (Anthropic), with build verification through Claude Code.
Computational search: \texttt{rh\_candidate\_search.py} (Python/NumPy).
Logan's son contributed the research ledger structure, the dagger-completion defect hypothesis,
and the FICS methodology framework.
ChatGPT and Claude are listed as tools, not authors.

\begin{thebibliography}{99}
\bibitem{BaezDuarte} L.~Báez-Duarte, A strengthening of the Nyman--Beurling criterion for the Riemann Hypothesis, \textit{Atti Accad.\ Naz.\ Lincei Cl.\ Sci.\ Fis.\ Mat.\ Natur.\ Rend.\ Lincei} \textbf{14} (2003), 5--11.
\bibitem{Beurling} A.~Beurling, A closure problem related to the Riemann zeta-function, \textit{Proc.\ Nat.\ Acad.\ Sci.\ USA} \textbf{41} (1955), 312--314.
\bibitem{Burnol} J.-F.~Burnol, Sur les formules explicites I: analyse invariante, \textit{C.~R.~Acad.\ Sci.\ Paris S\'{e}r.~I} \textbf{331} (2001), 423--428.
\bibitem{Connes} A.~Connes, Trace formula in noncommutative geometry and the zeros of the Riemann zeta function, \textit{Selecta Math.\ (N.S.)} \textbf{5} (1999), 29--106.
\bibitem{HongLoewy} S.~Hong and R.~Loewy, Asymptotic behavior of eigenvalues of GCD matrices, \textit{Glasg.\ Math.\ J.} \textbf{46} (2004), 551--569.
\bibitem{Mayer} D.~Mayer, The thermodynamic formalism approach to Selberg's zeta function for $\mathrm{PSL}(2,\mathbb{Z})$, \textit{Bull.\ Amer.\ Math.\ Soc.} \textbf{25} (1991), 55--60.
\bibitem{Montgomery} H.~L.~Montgomery, The pair correlation of zeros of the zeta function, in \textit{Analytic Number Theory}, Proc.\ Sympos.\ Pure Math.\ XXIV, AMS, 1973, 181--193.
\bibitem{Nyman} B.~Nyman, On the one-dimensional translation group and semi-group in certain function spaces, PhD thesis, Uppsala University, 1950.
\bibitem{RodgersTao} B.~Rodgers and T.~Tao, The de Bruijn--Newman constant is non-negative, \textit{Forum Math.\ Pi} \textbf{8} (2020), e6.
\bibitem{Weil} A.~Weil, Sur les ``formules explicites'' de la th\'{e}orie des nombres premiers, \textit{Comm.\ S\'{e}m.\ Math.\ Univ.\ Lund} (tome suppl\'{e}mentaire d\'{e}di\'{e} \`{a} Marcel Riesz), 1952, 252--265.
\bibitem{Mathlib} The Mathlib Community, Mathlib4, \url{https://github.com/leanprover-community/mathlib4}, 2023.
\end{thebibliography}

\end{document}

